La sirena d'una alarma d'un edifici emet una ona sonora de 0,136 m de longitud d'ona, que en l'aire es propaga a una velocitat d'uns 340 m/s. L'ona sonora arriba a un observador que està aturat en un semàfor, segons la direcció de l'eix $x$ en sentit positiu i amb una amplitud $A_0$.

\begin{apartat}{1,25}
Escriviu l'equació de l'ona harmònica plana, $A(x,t)$, que arriba a l'observador. Suposeu que la fase inicial és zero.
\end{apartat}

\begin{apartat}{1,25}
A continuació, el semàfor es posa verd, i l'observador es posa en moviment i s'acosta a l'alarma a una velocitat constant. En aquestes condicions, l'observador percep un canvi de la freqüència de l'ona. Quin canvi de la freqüència percep l'observador? Quin o quins dels arguments següents descriuen millor aquest fenomen?:

\begin{enumerate}
\item El canvi de la freqüència és degut al moviment de la font.
\item L'aparent canvi de la freqüència és degut al moviment relatiu entre la font i l'observador.
\item El canvi en la longitud d'ona de la font és degut al moviment de l'observador.
\item L'aparent canvi de la longitud d'ona és degut al moviment de l'observador.
\end{enumerate}

Justifiqueu la tria i per què heu descartat la resta d'arguments, i indiqueu el fenomen en què us heu basat.
\end{apartat}
